% The appendices that back the notes' claims: first how the claims were
% checked, then one chapter per claim that took a literature search.
% Each search chapter is written from that search's record and lives in
% literature-search/.

\chapter{How the claims in these notes were checked}
\label{app:checking}

These notes ask who checks, so they owe the reader the same answer
about themselves.
Much of the text was drafted with an AI agent, and none of its claims
was taken from what a language model remembers.
Each was checked in one of three ways.

\section{Claims about programs}
What a program does is shown by running it.
The output in \cref{ex:primes-run} is produced when these notes are
built, by the program printed in \cref{ex:primes}.
The conversations with the support bot are recordings of real runs,
with the model and the date in their headers, and the bot's own source,
with its tests, is printed in full in \cref{bot:chapter}.

\section{Claims that a source says something}
What a named source says (a news report, a standard, a court's
decision) was checked by opening that source and copying the passage.
The bibliography entry of every such source records the passage and the
date it was read.

\section{Claims about what is true in general}
A claim about how things are (who gains from AI assistance, whether an
attack can be prevented) cannot be settled by one source that agrees.
Each such claim was turned into a question and searched for in the
research literature in both directions: for studies that support it and
for studies that contradict or limit it.
The appendices that follow report those searches, one question each:
which databases were searched with which queries, what was found, and
what the notes may therefore say.

\section{Who did the checking}
The searches and the first reading of the sources were done by AI
agents, which recorded a verbatim passage for every source they cite.
An appendix marked \enquote{awaiting human review} has not yet been
gone through by the author.

% One chapter per backed claim, each written by its search.  A chapter
% is included once its search has also delivered the table of hits that
% the chapter ends with.
%
% A chapter's title is its research question, which is too long for the
% running head (two lines there make every page 12pt too tall), so each
% chapter gets a short form of the question for the head.
\newcommand{\searchhead}{}
\makeatletter
\renewcommand{\chaptermark}[1]{%
  \markboth{\@chapapp\ \thechapter.\ \searchhead}%
           {\@chapapp\ \thechapter.\ \searchhead}%
}
\makeatother
\newcommand{\searchchapter}[2]{%
  \renewcommand{\searchhead}{#2}%
  \IfFileExists{literature-search/#1-full.tex}%
    {\InputIfFileExists{literature-search/#1-chapter.tex}{}{}}%
    {}%
}
\searchchapter{who-gains}{Who gains more from AI assistance?}
\searchchapter{insecure-code}%
  {Is code written with an AI assistant less secure?}
\searchchapter{prompt-injection}%
  {Can prompt injection be reliably prevented?}
